\begin{problem}[理论考试][质点系与空气阻力]{31}
  如图，两个质量均为 $m$ 的小球 A 和 B（均可视为质点），固定在中心位于 C、长为 $2l$ 的刚性轻质细杆的两端，构成一质点系。在竖直面内建立 Oxy 坐标，Ox 方向沿水平向右，Oy 方向竖直向上。初始时质点系中心 C 位于原点 O，并以初速度 $v_{0}$ 竖直上抛，上抛过程中，ACB 三点连线始终水平。风速大小恒定为 $u$、方向沿 x 轴正向，小球在运动中所受空气阻力 $f$ 的大小与相对于空气运动速度 $v$ 的大小成正比，方向相反，即 $f = -kv$，$k$ 为正的常量。当 C 点升至最高点时，恰好有一沿 $y$ 轴正向运动、质量为 $m_1$、速度大小为 $u_1$ 的小石块（视为质点）与小球 A 发生竖直方向的碰撞，设碰撞是完全弹性的，碰撞时间极短。此后 C 点回落到上抛开始时的同一水平高度，此时它在 $Ox$ 方向上的位置记为 $s$，将从上抛到落回的整个过程所用时间记为 $T$，质点系旋转的圈数记为 $n$。求质点系
  \begin{subproblem}
    转动的初始角速度 $\omega_{0}$，及其回落到 s 点时角速度 $\omega_{s}$ 与 n 的关系；
  \end{subproblem}
  \begin{subproblem}
    从开始上抛到落回到 s 点为止的过程中，空气阻力做的功 $W_{f}$ 与 n、s、T 的关系。
  \end{subproblem}
\end{problem}